% !TeX program = pdflatex
\documentclass[11pt,a4paper]{amsart}

\newcommand{\notelanguage}{finnish}
% Shared preamble for course notes and exercise sets.

\usepackage[T1]{fontenc}
\usepackage{lmodern}
\providecommand{\notelanguage}{english}
\usepackage[finnish,english]{babel}
\AtBeginDocument{\expandafter\selectlanguage\expandafter{\notelanguage}}

\newcommand{\notesfinnishlanguage}{finnish}
\ifx\notelanguage\notesfinnishlanguage
\newcommand{\notestheoremname}{Lause}
\newcommand{\noteslemmaname}{Lemma}
\newcommand{\notespropositionname}{Propositio}
\newcommand{\notescorollaryname}{Korollaari}
\newcommand{\notesdefinitionname}{Määritelmä}
\newcommand{\notesexamplename}{Esimerkki}
\newcommand{\notesproblemname}{Tehtävä}
\newcommand{\notesexercisename}{Tehtävä}
\newcommand{\notesremarkname}{Huomautus}
\newcommand{\notessolutionname}{Ratkaisu}
\else
\newcommand{\notestheoremname}{Theorem}
\newcommand{\noteslemmaname}{Lemma}
\newcommand{\notespropositionname}{Proposition}
\newcommand{\notescorollaryname}{Corollary}
\newcommand{\notesdefinitionname}{Definition}
\newcommand{\notesexamplename}{Example}
\newcommand{\notesproblemname}{Problem}
\newcommand{\notesexercisename}{Exercise}
\newcommand{\notesremarkname}{Remark}
\newcommand{\notessolutionname}{Solution}
\fi

\usepackage[a4paper,margin=30mm]{geometry}
\usepackage{microtype}
\usepackage{graphicx}
\usepackage{enumitem}

\newcommand{\exerciseheading}[2]{%
  \par\noindent
  {\large\bfseries #1\par}
  \vspace{0.1em}
  \noindent{\large #2\par}
  \vspace{1.5em}
}

\usepackage{amsmath}
\usepackage{amssymb}
\usepackage{amsthm}
\usepackage{mathtools}
\usepackage{aliascnt}

% Row operations beside the affected row, e.g. \rowop{+2R_1}.
\newcommand{\rowop}[1]{\mathrel{\big|}\,#1}

% Linear systems: the argument gives the number of alignment pairs.
\newenvironment{linearsystem}[1]
{\left\{
  \begin{alignedat}{#1}}
    {
  \end{alignedat}\right.}

  \usepackage[hidelinks,pdfusetitle]{hyperref}

  \numberwithin{equation}{section}
  \setlist{itemsep=0.25em,topsep=0.5em}

  \theoremstyle{plain}
  \newtheorem{theorem}{\notestheoremname}[section]
  \newaliascnt{lemma}{theorem}
  \newtheorem{lemma}[lemma]{\noteslemmaname}
  \aliascntresetthe{lemma}
  \newaliascnt{proposition}{theorem}
  \newtheorem{proposition}[proposition]{\notespropositionname}
  \aliascntresetthe{proposition}
  \newaliascnt{corollary}{theorem}
  \newtheorem{corollary}[corollary]{\notescorollaryname}
  \aliascntresetthe{corollary}

  \theoremstyle{definition}
  \newaliascnt{definition}{theorem}
  \newtheorem{definition}[definition]{\notesdefinitionname}
  \aliascntresetthe{definition}
  \newaliascnt{example}{theorem}
  \newtheorem{example}[example]{\notesexamplename}
  \aliascntresetthe{example}
  \newaliascnt{problem}{theorem}
  \newtheorem{problem}[problem]{\notesproblemname}
  \aliascntresetthe{problem}
  \newtheorem{exercise}{\notesexercisename}

  \theoremstyle{remark}
  \newaliascnt{remark}{theorem}
  \newtheorem{remark}[remark]{\notesremarkname}
  \aliascntresetthe{remark}

  \newenvironment{solution}
  {
  \begin{proof}[\notessolutionname]}
    {
  \end{proof}}

  \usepackage[nameinlink,noabbrev]{cleveref}

  \crefname{theorem}{theorem}{theorems}
  \crefname{lemma}{lemma}{lemmas}
  \crefname{proposition}{proposition}{propositions}
  \crefname{corollary}{corollary}{corollaries}
  \crefname{definition}{definition}{definitions}
  \crefname{example}{example}{examples}
  \crefname{problem}{problem}{problems}
  \crefname{exercise}{exercise}{exercises}
  \crefname{remark}{remark}{remarks}

  \ifx\notelanguage\notesfinnishlanguage
  \crefname{theorem}{lause}{lauseet}
  \crefname{lemma}{lemma}{lemmat}
  \crefname{proposition}{propositio}{propositiot}
  \crefname{corollary}{korollaari}{korollaarit}
  \crefname{definition}{määritelmä}{määritelmät}
  \crefname{example}{esimerkki}{esimerkit}
  \crefname{problem}{tehtävä}{tehtävät}
  \crefname{exercise}{tehtävä}{tehtävät}
  \crefname{remark}{huomautus}{huomautukset}
  \fi

  \newcommand{\NN}{\mathbb{N}}
  \newcommand{\NNone}{\mathbb{N}_1}
  \newcommand{\ZZ}{\mathbb{Z}}
  \newcommand{\QQ}{\mathbb{Q}}
  \newcommand{\RR}{\mathbb{R}}
  \newcommand{\CC}{\mathbb{C}}
  \newcommand{\PP}{\mathbb{P}}

  \DeclareMathOperator{\im}{im}
  \DeclareMathOperator{\rank}{rank}
  \DeclareMathOperator{\syt}{syt}
  \DeclareMathOperator{\pyj}{pyj}
  \DeclareMathOperator{\dom}{dom}
  \DeclarePairedDelimiter{\abs}{\lvert}{\rvert}
  \DeclarePairedDelimiter{\norm}{\lVert}{\rVert}
  \newcommand{\set}[1]{\left\lbrace#1\right\rbrace}
  \newcommand{\maxset}[1]{\max\set{#1}}
  \newcommand{\minset}[1]{\min\set{#1}}
  \newcommand{\powerset}[1]{\mathcal{P}\!\left(#1\right)}


\title{Lukuteoria I\\Harjoitus 3}
\author{}
\date{}

\begin{document}

\exerciseheading{Lukuteoria I}{Harjoitus 3}

\begin{exercise}
  Henkilötunnuksen kuusi ensimmäistä numeroa ilmaisevat syntymäajan muodossa
  ppkkvv, ja välimerkki \texttt{-} tarkoittaa 1900-lukua. Tarkistusmerkki
  määräytyy syntymäajan ja kolminumeroisen yksilönumeron muodostaman
  yhdeksännumeroisen luvun jakojäännöksestä modulo \(31\); merkkiä
  \texttt{X} vastaa jakojäännös \(29\).
  Henkilötunnuksesta \texttt{2207x7-314X} on pyyhkiytynyt vuosiluvun
  ensimmäinen numero. Selvitä omistajan syntymävuosi.
\end{exercise}

\begin{exercise}
  \leavevmode\par
  \begin{enumerate}[label=(\alph*), nosep]
    \item Olkoon luvun \(n\) alkutekijähajotelma
      \(n = p_1^{\alpha_1}p_2^{\alpha_2}\cdots p_k^{\alpha_k}\).
      Osoita, että
      \[
        d(n) = \prod_{j=1}^{k}(\alpha_j + 1).
      \]
    \item Keksi esimerkki kokonaisluvusta \(n\), jolla on tasan \(20\)
      positiivista tekijää, ja laske \(\sigma(n)\).
  \end{enumerate}
\end{exercise}

\begin{exercise}
  \leavevmode\par
  \begin{enumerate}[label=(\alph*), nosep]
    \item Olkoot positiiviset kokonaisluvut \(m_1, m_2, \ldots, m_k\)
      pareittain yhteistekijättömiä eli \(\syt(m_i, m_j) = 1\) aina, kun
      \(i \ne j\). Olkoon
      \[
        x_0 = a_1M_1y_1 + a_2M_2y_2 + \cdots + a_kM_ky_k,
      \]
      missä \(M = m_1m_2\cdots m_k\), \(M_j = M/m_j\) ja \(y_j\)
      on luvun \(M_j\) käänteisluku modulo \(m_j\) eli
      \(y_jM_j \equiv 1 \pmod{m_j}\). Osoita, että \(x_0\) on
      kongruenssiyhtälöryhmän
      \[
        \begin{cases}
          x \equiv a_1 \pmod{m_1} \\
          x \equiv a_2 \pmod{m_2} \\
          \vdots \\
          x \equiv a_k \pmod{m_k}
        \end{cases}
      \]
      ratkaisu.
    \item Ratkaise kongruenssiyhtälöryhmä
      \[
        \begin{cases}
          x \equiv 3 \pmod{6} \\
          2x \equiv 3 \pmod{5} \\
          x \equiv -2 \pmod{7}.
        \end{cases}
      \]
  \end{enumerate}
\end{exercise}

\begin{exercise}
  Määritä luvun \(100!\) jakojäännös jaettaessa luvulla \(103\).
\end{exercise}

\begin{exercise}
  \leavevmode\par
  \begin{enumerate}[label=(\alph*), nosep]
    \item Tutki, ovatko luvut \(2016\), \(4675\) ja \(8128\) täydellisiä
      lukuja.
    \item Anna esimerkki sellaisesta parittomasta luvusta \(n\), että
      luentokalvojen (12.11.) lauseiden 4 ja 5 avulla ei voi suoraan
      todistaa, että \(n\) ei ole täydellinen luku.
  \end{enumerate}
\end{exercise}

\begin{exercise}
  Olkoon \(n\) pariton täydellinen luku. Osoita, että \(n\) ei ole
  jaollinen luvulla \(105\).

  \emph{Vinkki:} Luentomonisteen lauseesta 60 ja lausekkeen
  \(\sigma(n)/n\) arvioimisesta voisi olla hyötyä.
\end{exercise}

\end{document}
